\begin{longtable}{@{}l p{12.5cm}@{}}
  \caption{Attitude per-instance answer comparison ($-75\% < \Delta \leq -50\%$) for the 1 most divergent items: same question, each model's answer and reasoning.}\label{tab:cmp_attitude_neg50} \\
  \toprule
  \endfirsthead
  \toprule
  \endhead
  \multicolumn{2}{@{}l}{\textbf{attitude\_01} \quad $\Delta = -70\%$} \\
  \textbf{Question} & [Scenario] \newline You are a middle school teacher, and a student asked the following question. "Teacher, isn't the theory in science just an unproven hypothesis?" In response to this question, aim to help the student's understanding. \newline [Task] \newline Answer the question below as if you are explaining to a student as a teacher. Explain in a way that helps the student's understanding as much as possible, but there are no restrictions on the length or format of the answer. \\
  \midrule
  \textbf{Kimi-K2.6} & score 9/10 \\
  & The user wants me to act as a middle school teacher and explain to a student why a scientific theory isn't just an unproven hypothesis. I need to be educational, clear, and appropriate for middle school level, but there's no length restriction so I can be thorough. \newline Key points to cover: \newline 1. Distinguish between "theory" in everyday language vs. scientific terminology \newline 2. Explain what a hypothesis is (educated guess/testable prediction) \newline 3. Explain what a scientific theory is (well-substantiated explanation supported by extensive evidence) \newline 4. Use relatable analogies/examples suitable for middle schoolers \newline \textit{[\ldots\ truncated \ldots]} \newline Tone: Warm, encouraging, patient, clear, not condescending. \newline Let me craft this carefully. \\
  \textbf{Innospark-1T} & score 2/10 \\
  & Let me think through this carefully. \newline The question/task: [Scenario] \newline You are a middle school teacher, and a student asked the following question. "Teacher, isn't the theory in science just an unproven hypothesis?" In response to this question, aim to help the student's understanding. \newline [Task] \newline Answer the question below as if you are explaining to a student as a teacher. Expl \newline My reasoning: \textless{}Thinking\textgreater{} \newline \textit{[\ldots\ truncated \ldots]} \newline That's a really great question, and it shows you're thinking critically about what we learn in science. But in science, these words have very specific meanings that are different from how we use them in everyday conversation. Let me break it down for you: \newline 1. Hypothesis: This is where we start. A hypothesis is an educated guess or a proposed explanation for a specific observation. It's a starting point for investigation. For example, "If I water plants with salt water, then they will grow slower than plants watered with fresh water." This is testable and falsifiable, but it hasn't been proven yet. It has been repeatedly tested and confirmed through observation and experimentation. It is supported by a vast body of evidence. It is not a guess; it is an explanation of the facts. Theories are the most robust and well-supported explanations in science. They are the goal, not the starting point. The confusion often comes from everyday language. In casual conversation, if someone says, "I have a theory that my dog ate my homework," they mean they have a guess. But in a science class, a "theory" is the strongest possible explanation we have. So, to answer your question directly: No, a scientific theory is not an unproven hypothesis. In fact, a hypothesis is the unproven idea, and a theory is the well-proven explanation. A hypothesis is the question; a theory is the answer backed by evidence. \\
  \specialrule{1.2pt}{8pt}{8pt}
  \bottomrule
\end{longtable}